% ======================================================================
\chapter[Атомное ядро. Радиус ядра. Формула Вайцзеккера]%
        {Атомное ядро. Радиус ядра. Формула Вайцзеккера}
\label{t:17}
% ======================================================================

\begin{qbox}
\textbf{Билет 17.} Атомное ядро, зависимость радиуса от массового числа. Формула Вайцзеккера
для энергии связи ядра.
\end{qbox}

\section{Содержание билета}

\textbf{Подтемы:}
\begin{enumerate}
 \item состав ядра, обозначения, изотопы/изобары/изотоны; свойства ядерных сил;
 \item размеры ядер: методы измерения, закон $R=r_0A^{1/3}$, распределение плотности, постоянство плотности;
 \item дефект массы и энергия связи; удельная энергия связи и её зависимость от $A$;
 \item капельная модель и полуэмпирическая формула Вайцзеккера: смысл и обоснование каждого члена;
 \item следствия: долина стабильности, энергетика деления и синтеза.
\end{enumerate}

\textbf{Необходимые понятия:} $E=mc^2$; электростатическая энергия заряженного шара; газ Ферми
(\tref{t:04}: число состояний в фазовом объёме; \tref{t:13}: принцип Паули); сильное взаимодействие (\tref{t:15}).

\section{Определения и физические принципы}

\begin{definition}[ядро]
Атомное ядро — связанная система $Z$ протонов и $N$ нейтронов (нуклонов); $A=Z+N$ — массовое число.
Обозначение $\nuc{X}{A}{Z}$ (например, $\nuc{U}{238}{92}$). \emph{Изотопы} — одинаковые $Z$, \emph{изобары} — одинаковые $A$,
\emph{изотоны} — одинаковые $N$.
\end{definition}

\begin{definition}[энергия связи]
Энергия, необходимая для разделения ядра на свободные нуклоны:
\begin{equation}\label{eq:17:Eb}
  E_{\text{св}}(A,Z)=\bigl[Zm_p+Nm_n-M_{\text{я}}(A,Z)\bigr]c^2=\bigl[ZM(\nucl{H}{1})+Nm_n-M_{\text{ат}}(A,Z)\bigr]c^2,
\end{equation}
где во втором выражении использованы массы нейтральных атомов (массы электронов сокращаются; разностью
энергий связи электронов пренебрегаем). Величина $\Delta M=E_{\text{св}}/c^2$ — \emph{дефект массы};
$\varepsilon=E_{\text{св}}/A$ — \emph{удельная энергия связи}.
\end{definition}

\begin{definition}[свойства ядерных сил]
Ядерные силы — сильное (остаточное цветовое) взаимодействие нуклонов: (1) короткодействующие, радиус
$\sim1{,}5$--$2$~фм; (2) притягивающие на расстояниях $\sim1$~фм и отталкивающие на расстояниях $\lesssim0{,}5$~фм (кор);
(3) зарядово-независимые ($pp\approx nn\approx np$); (4) спин-зависимые и нецентральные; (5) обладают свойством
\emph{насыщения}: нуклон взаимодействует лишь с ближайшими соседями.
\end{definition}

\emph{Область применимости.} Формула Вайцзеккера — полуэмпирическая (капельная модель); она хорошо описывает
энергии связи средних и тяжёлых ядер (точность порядка 1\,\% и лучше), но не воспроизводит оболочечные эффекты
(магические числа 2, 8, 20, 28, 50, 82, 126) и неприменима к самым лёгким ядрам.

\section{Основная теория}

\subsection{Размеры ядер}

\paragraph{Методы.}
\begin{enumerate}
 \item \textbf{Рассеяние $\alpha$-частиц} (Резерфорд): отклонения от формулы Резерфорда при больших энергиях и углах
       начинаются, когда частица подходит к ядру на расстояние действия ядерных сил.
 \item \textbf{Рассеяние быстрых электронов} (Хофштадтер, 1950-е; Нобелевская премия 1961): электроны с $\lambda\sim1$~фм
       (энергии сотни МэВ) не чувствуют ядерных сил; дифракционная картина даёт распределение \emph{заряда} в ядре.
 \item \textbf{Мюонные атомы} (\tref{t:03}): уровни мюона, проводящего часть времени внутри ядра, сдвинуты относительно
       точечного ядра.
 \item \textbf{Зеркальные ядра} (вывод~\ref{der:17:mirror}): разность энергий связи пар ($\nucl{N}{15}$--$\nucl{O}{15}$ и др.) обусловлена
       кулоновской энергией и даёт $R$.
 \item \textbf{Рассеяние нейтронов} и времена жизни $\alpha$-распада — радиус действия ядерных сил.
\end{enumerate}

\paragraph{Результат.} Все методы дают
\begin{equation}\label{eq:17:R}
  R=r_0A^{1/3},\qquad r_0\approx1{,}2\ \text{фм (по распределению заряда)},\quad 1{,}3\text{--}1{,}4\ \text{фм (по ядерным силам)}.
\end{equation}
Различие связано с тем, что ядерные силы действуют на некотором расстоянии за пределами распределения
нуклонов. Объём ядра пропорционален числу нуклонов: $V=\frac43\pi r_0^3A$, т.\,е. \emph{плотность ядерного вещества
одинакова для всех ядер} (вывод~\ref{der:17:density}):
\[
  n\approx\frac{3}{4\pi r_0^3}\approx0{,}14\text{--}0{,}17\ \text{фм}^{-3},\qquad \rho\approx2{,}3\cdot10^{14}\ \text{г/см}^3.
\]
Это прямое следствие насыщения ядерных сил и отталкивания на малых расстояниях: нуклоны «упакованы» как молекулы
в капле жидкости.

\paragraph{Распределение плотности} (из рассеяния электронов) описывается распределением Ферми
\begin{equation}\label{eq:17:fermi}
  \rho(r)=\frac{\rho_0}{1+\exp\bigl((r-c)/a\bigr)},\qquad c\approx1{,}07A^{1/3}\ \text{фм},\quad a\approx0{,}55\ \text{фм}:
\end{equation}
почти постоянная плотность в центре и размытый поверхностный слой толщиной $t\approx4{,}4a\approx2{,}4$~фм (спад
от 90 до 10\,\% $\rho_0$) — одинаковой для всех ядер (рис.~\ref{fig:17:fermi}).

\begin{figure}[htbp]
\centering
\begin{tikzpicture}
\begin{axis}[width=0.64\linewidth, height=4.6cm, xlabel={$r$, фм}, ylabel={$\rho/\rho_0$}, xmin=0, xmax=11, ymin=0, ymax=1.15,
   samples=200, legend style={font=\small, draw=none}, legend pos=north east]
  \addplot[very thick, accent, domain=0:11] {1/(1+exp((x-1.07*16^(1/3))/0.55))};
  \addlegendentry{$^{16}$O}
  \addplot[very thick, keycol, domain=0:11] {1/(1+exp((x-1.07*208^(1/3))/0.55))};
  \addlegendentry{$^{208}$Pb}
\end{axis}
\end{tikzpicture}
\caption{Распределение Ферми~\eqref{eq:17:fermi} для лёгкого и тяжёлого ядер (расчёт по формуле с $c=1{,}07A^{1/3}$~фм,
$a=0{,}55$~фм). Центральная плотность одинакова, толщина поверхностного слоя одинакова.}
\label{fig:17:fermi}
\end{figure}

\subsection{Энергия связи. Удельная энергия связи}

Примеры: дейтрон — $2{,}22$~МэВ ($1{,}11$~МэВ/нуклон); $\nucl{He}{4}$ — $28{,}3$~МэВ ($7{,}07$~МэВ/нуклон); $\nucl{Fe}{56}$ — $8{,}79$~МэВ/нуклон;
$\nucl{U}{238}$ — $7{,}57$~МэВ/нуклон. Основные закономерности:
\begin{itemize}
 \item для $A\gtrsim20$ удельная энергия связи почти постоянна, $\varepsilon\approx7{,}5$--$8{,}8$~МэВ, т.\,е. $E_{\text{св}}\propto A$ —
       ещё одно проявление насыщения (если бы каждый нуклон взаимодействовал со всеми, было бы $E_{\text{св}}\propto A^2$);
 \item максимум $\varepsilon$ в области железа--никеля ($A\approx56$--$62$); слева — рост, справа — плавный спад;
 \item у лёгких ядер $\varepsilon$ меньше и испытывает скачки (особо устойчивы $\nucl{He}{4}$, $\nucl{C}{12}$, $\nucl{O}{16}$).
\end{itemize}
Энергия связи — около 1\,\% энергии покоя ядра ($8\MeV/939\MeV$), тогда как в атоме — $\sim10^{-8}$.

\subsection{Капельная модель и формула Вайцзеккера}

Постоянство плотности и насыщение сил делают ядро похожим на каплю несжимаемой заряженной жидкости. Энергия связи
складывается из вкладов, каждый из которых имеет ясную физическую природу (К.~Вайцзеккер, 1935):
\begin{equation}\label{eq:17:SEMF}
  \boxed{E_{\text{св}}(A,Z)=a_VA-a_SA^{2/3}-a_C\frac{Z^2}{A^{1/3}}-a_A\frac{(A-2Z)^2}{A}+\delta(A,Z)}
\end{equation}
\begin{enumerate}
 \item \textbf{Объёмная энергия} $a_VA$: каждый нуклон связан с фиксированным числом соседей (насыщение) и даёт
       одинаковый вклад; $E\propto V\propto A$.
 \item \textbf{Поверхностная энергия} $-a_SA^{2/3}$: нуклоны на поверхности имеют меньше соседей и связаны слабее; поправка
       пропорциональна площади $4\pi R^2\propto A^{2/3}$ (аналог поверхностного натяжения). Относительно наиболее важна для лёгких ядер.
 \item \textbf{Кулоновская энергия} $-a_CZ^2/A^{1/3}$: электростатическое отталкивание протонов, равномерно распределённых по шару
       радиуса $R=r_0A^{1/3}$: $E_C=\frac35\frac{Z^2e^2}{R}$ (вывод~\ref{der:17:coulomb}); $a_C=\frac{3e^2}{5r_0}\approx0{,}7$~МэВ. (Иногда пишут
       $Z(Z-1)$, исключая самодействие протона.) Определяет спад $\varepsilon$ у тяжёлых ядер.
 \item \textbf{Энергия симметрии (асимметрии)} $-a_A(N-Z)^2/A$: при заданном $A$ энергетически выгодно $N=Z$. Причина —
       принцип Паули: протоны и нейтроны заполняют отдельные системы уровней, и избыток частиц одного сорта
       вынужден занимать более высокие уровни (вывод~\ref{der:17:asym}). Форма $(N-Z)^2/A$ следует из модели газа Ферми.
 \item \textbf{Энергия спаривания} $\delta$: нуклоны одного сорта образуют пары с противоположными спинами, что дополнительно
       выгодно:
       \[
         \delta=\begin{cases}+a_PA^{-3/4}, & Z,N\ \text{чётные (чётно-чётные ядра)},\\ 0, & A\ \text{нечётное},\\ -a_PA^{-3/4}, & Z,N\ \text{нечётные}.\end{cases}
       \]
       Поэтому стабильные нечётно-нечётные ядра почти отсутствуют (лишь $\nucl{H}{2}$, $\nucl{Li}{6}$, $\nucl{B}{10}$, $\nucl{N}{14}$).
\end{enumerate}
Один из распространённых наборов коэффициентов (подбираются по экспериментальным массам; разные авторы дают немного
различающиеся значения):
\begin{equation}\label{eq:17:coef}
  a_V=15{,}75,\quad a_S=17{,}8,\quad a_C=0{,}711,\quad a_A=23{,}7,\quad a_P=34\ \ \text{(МэВ)}.
\end{equation}
Проверка для $\nucl{Fe}{56}$ и $\nucl{U}{238}$ — в выводе~\ref{der:17:check}. Вклады членов в $\varepsilon(A)$ показаны на рис.~\ref{fig:17:BA}.

\begin{figure}[htbp]
\centering
\begin{tikzpicture}
\begin{axis}[width=0.86\linewidth, height=6.6cm, xlabel={$A$}, ylabel={$E_{\text{св}}/A$, МэВ}, xmin=0, xmax=250, ymin=0, ymax=17.5,
   samples=120, legend style={font=\small, draw=none, fill=white, fill opacity=0.85, text opacity=1}, legend pos=south east,
   declare function={Z0(\A)=\A/(1.98+0.0149*\A^(2/3));}]
  \addplot[gray, thick, dashed, domain=4:250] {15.75};
  \addlegendentry{объёмный член}
  \addplot[dercol, thick, domain=4:250] {15.75-17.8*x^(-1/3)};
  \addlegendentry{$-$ поверхность}
  \addplot[keycol, thick, domain=4:250] {15.75-17.8*x^(-1/3)-0.711*Z0(x)^2/x^(4/3)};
  \addlegendentry{$-$ кулоновский}
  \addplot[accent, very thick, domain=4:250] {15.75-17.8*x^(-1/3)-0.711*Z0(x)^2/x^(4/3)-23.7*(x-2*Z0(x))^2/x^2};
  \addlegendentry{$-$ симметрия (итог)}
\end{axis}
\end{tikzpicture}
\caption{Удельная энергия связи по формуле Вайцзеккера (без члена спаривания) вдоль долины стабильности $Z=Z_0(A)$
и вклады отдельных членов (расчёт по коэффициентам~\eqref{eq:17:coef}). Максимум $\approx8{,}8$~МэВ при $A\approx60$;
для лёгких ядер формула неточна.}
\label{fig:17:BA}
\end{figure}

\subsection{Следствия формулы Вайцзеккера}

\paragraph{Долина стабильности.} Масса ядра при фиксированном $A$ — парабола по $Z$ (вывод~\ref{der:17:valley}); её минимум
\begin{equation}\label{eq:17:Z0}
  Z_0(A)\approx\frac{A}{1{,}98+0{,}015A^{2/3}}.
\end{equation}
Для лёгких ядер $Z_0\approx A/2$ ($N\approx Z$), для тяжёлых кулоновское отталкивание сдвигает стабильность к избытку
нейтронов ($\nucl{U}{238}$: $Z_0\approx93$, реально $Z=92$, $N/Z\approx1{,}6$). Ядра выше и ниже долины испытывают бета-распады
(\tref{t:20}).

\paragraph{Энергетика деления и синтеза.} Из хода $\varepsilon(A)$: при делении тяжёлого ядра ($\varepsilon\approx7{,}6$~МэВ) на два средних
($\varepsilon\approx8{,}5$~МэВ) выделяется $\approx0{,}9\cdot236\approx200$~МэВ; при синтезе лёгких ядер ($\nucl{H}{2}+\nucl{H}{3}\to\nucl{He}{4}+n$) —
$17{,}6$~МэВ, т.\,е. больше на нуклон. Деление энергетически выгодно, когда кулоновская энергия велика по сравнению с
поверхностной; параметр делимости $Z^2/A$: при $Z^2/A\gtrsim2a_S/a_C\approx50$ ядро неустойчиво к делению мгновенно.

\paragraph{Альфа-распад} энергетически возможен для $A\gtrsim150$ (\tref{t:19}): энергия связи $\alpha$-частицы ($28{,}3$~МэВ)
превышает $4\varepsilon$ тяжёлого ядра.

\section{Математические выводы}

\begin{derivation}[label=der:17:density]{плотность ядерного вещества}
Число нуклонов в единице объёма:
\[
  n=\frac{A}{\frac43\pi R^3}=\frac{A}{\frac43\pi r_0^3A}=\frac{3}{4\pi r_0^3}=\frac{3}{4\pi(1{,}2\ \text{фм})^3}\approx0{,}14\ \text{фм}^{-3}
\]
— от $A$ не зависит. Массовая плотность $\rho=nm_N=0{,}14\cdot10^{39}\ \text{см}^{-3}\cdot1{,}67\cdot10^{-24}\ \text{г}\approx2{,}3\cdot10^{14}$~г/см$^3$
(центральная плотность из~\eqref{eq:17:fermi} немного выше, $\approx0{,}16$~фм$^{-3}$). Среднее расстояние между нуклонами
$n^{-1/3}\approx1{,}9$~фм — порядка радиуса ядерных сил.
\end{derivation}

\begin{derivation}[label=der:17:coulomb]{кулоновская энергия равномерно заряженного шара}
Соберём шар радиуса $R$ с зарядом $Q$ послойно. Когда уже собран шар радиуса $r$ с зарядом $q(r)=Q r^3/R^3$, доставка из
бесконечности слоя $\dd q=3Qr^2\dd r/R^3$ требует работы $\dd W=q(r)\,\dd q/r$. Тогда
\[
  E_C=\int_0^R\frac{q(r)}{r}\dd q=\int_0^R\frac{Qr^3}{R^3r}\cdot\frac{3Qr^2}{R^3}\dd r=\frac{3Q^2}{R^6}\int_0^Rr^4\dd r=\frac35\frac{Q^2}{R}.
\]
При $Q=Ze$, $R=r_0A^{1/3}$: $E_C=\dfrac{3e^2}{5r_0}\dfrac{Z^2}{A^{1/3}}$, т.\,е. $a_C=\dfrac{3}{5}\dfrac{e^2}{r_0}=0{,}6\cdot\dfrac{1{,}44\ \text{МэВ}\cdot\text{фм}}{1{,}2\ \text{фм}}=0{,}72$~МэВ —
в согласии с эмпирическим $0{,}71$~МэВ.
\end{derivation}

\begin{derivation}[label=der:17:asym]{энергия симметрии в модели газа Ферми}
\emph{Допущения:} нуклоны — свободные фермионы со спином $\frac12$ в объёме $V$ ядра; протоны и нейтроны — две независимые
системы; учитывается только кинетическая энергия.

\emph{Шаг 1.} Число состояний с импульсом меньше $p_F$: $2\cdot V\cdot\frac43\pi p_F^3/(2\pi\hbar)^3$ (множитель 2 — спин). Приравняв числу
частиц $N$: $p_F=\hbar(3\pi^2N/V)^{1/3}$, $\varepsilon_F=p_F^2/2m$.

\emph{Шаг 2.} Средняя кинетическая энергия частицы $\frac35\varepsilon_F$, полная $E=\frac35N\varepsilon_F\propto N^{5/3}/V^{2/3}$. Для двух сортов:
\[
  E_{\text{кин}}=C\,\frac{N^{5/3}+Z^{5/3}}{A^{2/3}},\qquad C=\frac35\frac{\hbar^2}{2m}\Bigl(\frac{3\pi^2}{\frac43\pi r_0^3}\Bigr)^{2/3}\quad(V\propto A).
\]
\emph{Шаг 3.} Положим $N=\frac A2(1+\eta)$, $Z=\frac A2(1-\eta)$, $\eta=(N-Z)/A$. Разложение до $\eta^2$:
$(1\pm\eta)^{5/3}\approx1\pm\frac53\eta+\frac59\eta^2$, поэтому
\[
  N^{5/3}+Z^{5/3}=\Bigl(\frac A2\Bigr)^{5/3}\Bigl(2+\frac{10}{9}\eta^2\Bigr).
\]
\emph{Шаг 4.} Отсюда $E_{\text{кин}}=E_0\bigl(1+\frac59\eta^2\bigr)$, где $E_0=\frac35A\varepsilon_F^{(0)}$ — энергия симметричного ядра ($\varepsilon_F^{(0)}$ —
энергия Ферми при $N=Z$). Добавка
\[
  \Delta E=\frac35A\varepsilon_F^{(0)}\cdot\frac59\eta^2=\frac{\varepsilon_F^{(0)}}{3}\frac{(N-Z)^2}{A}.
\]
При $\varepsilon_F^{(0)}\approx37$~МэВ кинетический вклад $\approx12$~МэВ — около половины эмпирического $a_A$; остальное даёт зависимость
ядерного взаимодействия от изоспина ($np$-притяжение сильнее). Главное — модель объясняет \emph{форму} $(N-Z)^2/A$.
\end{derivation}

\begin{derivation}[label=der:17:valley]{долина стабильности}
Масса атома $M(A,Z)c^2=Zm_{\mathrm H}c^2+(A-Z)m_nc^2-E_{\text{св}}(A,Z)$ (с массой атома водорода). При фиксированном $A$:
\[
  \pdiff{M}{Z}c^2=(m_{\mathrm H}-m_n)c^2+\frac{2a_CZ}{A^{1/3}}-\frac{4a_A(A-2Z)}{A}=0 .
\]
Решая относительно $Z$, с $(m_n-m_{\mathrm H})c^2=0{,}78$~МэВ:
\[
  Z_0=\frac{4a_A+0{,}78}{\dfrac{8a_A}{A}+\dfrac{2a_C}{A^{1/3}}}=\frac{A}{\dfrac{8a_A}{4a_A+0{,}78}+\dfrac{2a_C}{4a_A+0{,}78}A^{2/3}}
  =\frac{A}{1{,}98+0{,}0149A^{2/3}}.
\]
Поскольку $M$ квадратична по $Z$, изобары лежат на параболе $M(Z)=M(Z_0)+\frac12\bigl(\frac{8a_A}{A}+\frac{2a_C}{A^{1/3}}\bigr)(Z-Z_0)^2/c^2$
(плюс член спаривания).
\end{derivation}

\begin{derivation}[label=der:17:mirror]{радиус по зеркальным ядрам}
Зеркальные ядра (например, $\nucl{N}{15}$: $Z=7$, $N=8$ и $\nucl{O}{15}$: $Z=8$, $N=7$) отличаются заменой $p\leftrightarrow n$. Из-за
зарядовой независимости ядерные энергии одинаковы, и разность энергий связи равна разности кулоновских:
\[
  \Delta E_C=\frac35\frac{e^2}{R}\bigl[Z^2-(Z-1)^2\bigr]=\frac35\frac{e^2}{R}(2Z-1).
\]
Измерив $\Delta E_C$ (по энергии $\beta^+$-распада $Z\to Z-1$ с учётом $m_n-m_p$), находим $R$. Результат согласуется с $r_0\approx1{,}2$--$1{,}3$~фм.
\end{derivation}

\begin{derivation}[label=der:17:check]{численная проверка формулы}
\emph{$\nucl{Fe}{56}$} ($Z=26$, $A^{1/3}=3{,}826$, $A^{2/3}=14{,}64$, чётно-чётное):
$15{,}75\cdot56=882{,}0$; $17{,}8\cdot14{,}64=260{,}6$; $0{,}711\cdot676/3{,}826=125{,}6$; $23{,}7\cdot16/56=6{,}8$; $\delta=34/56^{3/4}=1{,}7$.
$E_{\text{св}}\approx882{,}0-260{,}6-125{,}6-6{,}8+1{,}7=490{,}7$~МэВ, $\varepsilon\approx8{,}76$~МэВ (опыт: 492,3~МэВ; 8,79).

\emph{$\nucl{U}{238}$} ($Z=92$, $A^{1/3}=6{,}197$, $A^{2/3}=38{,}40$): $3748{,}5-683{,}5-971{,}1-290{,}4+0{,}6=1804$~МэВ, $\varepsilon\approx7{,}58$~МэВ
(опыт: 7,57). Пять параметров описывают энергии связи сотен ядер с точностью $\sim0{,}5$--$1\,\%$.
\end{derivation}

\section{Итоговые формулы}

\keyformula{R=r_0A^{1/3},\qquad r_0\approx1{,}2\text{--}1{,}4\ \text{фм},\qquad n=\frac{3}{4\pi r_0^3}\approx0{,}15\ \text{фм}^{-3}}
 {$A$ — массовое число; $n$ — концентрация нуклонов.}
 {все ядра, кроме самых лёгких; $r_0$ зависит от метода (заряд или ядерные силы).}
 {размер ядра; постоянство ядерной плотности.}

\keyformula{E_{\text{св}}=\bigl[Zm_{\mathrm H}+Nm_n-M_{\text{ат}}\bigr]c^2}
 {$m_{\mathrm H}$ — масса атома водорода, $M_{\text{ат}}$ — масса атома.}
 {точно (с пренебрежением разностью электронных энергий связи).}
 {энергия связи по таблицам масс; энерговыделение реакций.}

\keyformula{E_{\text{св}}=a_VA-a_SA^{2/3}-a_C\frac{Z^2}{A^{1/3}}-a_A\frac{(A-2Z)^2}{A}+\delta}
 {$a_V\approx15{,}8$, $a_S\approx17{,}8$, $a_C\approx0{,}71$, $a_A\approx23{,}7$~МэВ; $\delta=\pm34A^{-3/4}$~МэВ или 0.}
 {средние и тяжёлые ядра; без оболочечных поправок.}
 {энергии связи, массы, долина стабильности, энергии $\alpha$- и $\beta$-распадов, деления.}

\keyformula{a_C=\frac35\frac{e^2}{r_0},\qquad Z_0\approx\frac{A}{1{,}98+0{,}015A^{2/3}}}
 {$e^2=1{,}44$~МэВ$\cdot$фм; $Z_0$ — заряд наиболее стабильного изобара.}
 {капельная модель.}
 {связь коэффициента с радиусом; положение стабильных ядер.}

\section{Физическая интерпретация}

\textbf{Насыщение.} Пропорциональность $V\propto A$ и $E_{\text{св}}\propto A$ — два независимых проявления того, что ядерные силы
короткодействующие и насыщаются; без отталкивания на малых расстояниях ядро сжималось бы с ростом $A$.

\textbf{Конкуренция членов.} Поверхностный член уменьшает $\varepsilon$ у лёгких ядер, кулоновский — у тяжёлых; отсюда максимум
около железа. Это определяет и источники ядерной энергии (деление тяжёлых, синтез лёгких), и конечный продукт
звёздного нуклеосинтеза (железный пик). Существование тяжёлых ядер ограничено кулоновским отталкиванием (деление,
$\alpha$-распад).

\textbf{Ограничения модели.} Отклонения от формулы Вайцзеккера на несколько МэВ вблизи магических чисел указывают на
оболочечную структуру ядра — нуклоны движутся в общем потенциале почти независимо, как электроны в атоме.

\section{Что необходимо знать на экзамене}

\begin{examcheck}
\checkhead{Определения}
\begin{itemize}
 \item состав ядра, изотопы/изобары/изотоны; свойства ядерных сил (насыщение!);
 \item дефект массы, энергия связи, удельная энергия связи.
\end{itemize}
\checkhead{Формулы}
\begin{itemize}
 \item $R=r_0A^{1/3}$, $r_0\approx1{,}2$~фм; $n\approx0{,}15$~фм$^{-3}$; распределение Ферми (качественно);
 \item формула Вайцзеккера со всеми пятью членами и порядками коэффициентов; $Z_0(A)$.
\end{itemize}
\checkhead{Выводы}
\begin{itemize}
 \item постоянство плотности из $R\propto A^{1/3}$; кулоновская энергия шара $\frac35Q^2/R$ и оценка $a_C$;
 \item форма энергии симметрии из газа Ферми; долина стабильности; (желательно) метод зеркальных ядер.
\end{itemize}
\checkhead{Интерпретации}
\begin{itemize}
 \item физический смысл каждого члена; ход кривой $\varepsilon(A)$ и максимум у Fe;
 \item энергия деления ($\sim200$~МэВ) и синтеза; почему тяжёлые ядра нейтроноизбыточны; методы измерения $R$.
\end{itemize}
\end{examcheck}
